\scanpage{78}
% archive foliation: 77

\uline{Conclusion} Les catégories $\uline{\operatorname{Eq}}(C)$,
$C\in\operatorname{Ob}\underline B_0$, sont
\deleted{des \unsure{gr-catégories munies de structures}}
\inserted{munies canoniquement d'une structure de}
catégorie de Picard \uline{strictes}
\deleted{\unsure{commutatives}}.
\deleted{\unsure{Ils sont}}
Ces cat. de Picard sont \uline{canoniquement} équivalentes entre elles,
et peuvent donc s'identifier à une même catégorie de Picard
$\underline B$ (catégorie de Picard de \uline{Brandt}).
Donc $\underline B$ opère sur chaque $C$, et pour
$\unsure{C''|C},C'\in\operatorname{Ob}\underline B_0$,
la catégorie \deleted{\unsure{Hom}}
$\uline{\operatorname{Eq}}(C,C')$ est canoniquement munie d'une structure
de tors-catégorie sous $\underline B$. De façon précise,
\[
C'\longmapsto\uline{\operatorname{Eq}}(C,C')
\]
\[
\underline B_0\longrightarrow
\text{la catégorie des tors-catégories sous }\underline B
\]
est une équivalence de $2$-catégories.

Si on a un chgt de base
\diagram{
\begin{array}{ccc}
A&\longrightarrow&A'\\
\unsure{\cap}&&\unsure{\cap}\\
K&\unsure{\subset|-}&K'
\end{array}
}
et si les $C_{A'}$ satisfont les mêmes conditions que les $C$,
alors on a un hom de \deleted{\unsure{gr}} catégories de Picard
\[
\underline B\longrightarrow\underline B'
\]
et l'hom \deleted{\unsure{de base}}
\[
C\longrightarrow C_{A'}
\]
est compatible avec les opérations de $\underline B,\underline B'$ \ldots
